% Part: incompleteness
% Chapter: representability-in-q
% Section: basic-representable

\documentclass[../../../include/open-logic-section]{subfiles}

\begin{document}

\olfileid{inc}{req}{bre}
\olsection{بنسټيزې تابعې په~$\Th{Q}$ کښې تمثيلېدونکې دي}

لومړے بايد وښيو چې ټولې بنسټيزې تابعې په~$\Th{Q}$ کښې
تمثيلېدونکې دي۔ په پاې کښې بايد وښيو چې هرې~$k$-ځايه بنسټيزې
تابعې~$f(x_0,\dots,x_{k-1})$ ته داسې !!a{formula}
$!A_f(x_0,\dots,x_{k-1},y)$ څنګه وټاکو چې هغه تمثيل کړي۔

صفر، تالي، جمع، ضرب، د مساواتو ځانګړونکې تابعه، او پروجکشنونه
تمثيلولے شو۔ په هره قضيه کښې مناسب تمثيلوونکی فارمول په بشپړه توګه
ساده دے؛ د بېلګې په توګه، صفر د~$y = \Obj 0$ فارمول په وسيله
تمثيلېږي، تالي د~$x_0' = y$ !!{formula} په وسيله تمثيلېږي، او جمع
د~$(x_0 + x_1) = y$ !!{formula} په وسيله تمثيلېږي۔ کار په دې کښې دے چې وښيو
$\Th{Q}$ اړوند !!{sentence}s ثابتولے شي؛ د بېلګې په توګه، دا ويل
چې جمع د پورته !!{formula} په وسيله تمثيل شوې ده، دا غواړي چې وښيو
د طبيعي عددونو د هرې جوړې~$m$ او~$n$ دپاره~$\Th{Q}$ دا ثابتوي:
\begin{align*}
& \eq[\num n + \num m][\num {n+m}] \text{ او}\\
& \lforall[y][(\eq[(\num n + \num m)][y] \lif \eq[y][\num{n+m}])].
\end{align*}

\begin{prop}
\ollabel{prop:rep-zero}
صفر تابعه~$\Zero(x) = 0$ په~$\Th{Q}$ کښې د
$!A_{\Zero}(x,y) \ident \eq[y][\Obj 0]$ په وسيله تمثيلېږي۔
\end{prop}

\begin{prop}
\ollabel{prop:rep-succ}
تالي تابعه~$\Succ(x) = x+1$ په~$\Th{Q}$ کښې د
$!A_{\Succ}(x,y) \ident \eq[y][x']$ په وسيله تمثيلېږي۔
\end{prop}

\begin{prop}
\ollabel{prop:rep-proj}
پروجکشن تابعه~$\Proj{n}{i}(x_0, \dots, x_{n-1}) = x_i$ په
$\Th{Q}$ کښې د دې په وسيله تمثيلېږي:
\[!A_{\Proj{n}{i}}(x_0, \dots, x_{n-1}, y) \ident \eq[y][x_i].\]
\end{prop}

\begin{prob}
ثابته کړئ چې~$\eq[y][\Obj 0]$، $\eq[y][x']$، او~$\eq[y][x_i]$ په
ترتيب سره~$\Zero$، $\Succ$، او~$\Proj{n}{i}$ تمثيلوي۔
\end{prob}

\begin{prop}
\ollabel{prop:rep-id}
د~$=$ ځانګړونکې تابعه،
\[
\Char{=}(x_0, x_1) =
\begin{cases}
  1 & \text{که } x_0 =x_1\\
  0 & \text{بل ډول}
\end{cases}
\]
په~$\Th{Q}$ کښې د دې په وسيله تمثيلېږي:
\[
  !A_{\Char{=}}(x_0, x_1, y) \ident (\eq[x_0][x_1] \land \eq[y][\num{1}]) \lor (\eq/[x_0][x_1] \land
\eq[y][\num{0}]).
\]
\end{prop}

ثبوت لاندې لمې ته اړتيا لري۔

\begin{lem}
\ollabel{lem:q-proves-neq} که طبيعي عددونه~$n$ او~$m$ داسې وي چې
$n \neq m$، نو $\Th{Q} \Proves \eq/[\num n][\num m]$۔
\end{lem}

\begin{proof}
پر~$n$ استقرا وکاروئ څو وښيئ چې د هر~$m$ دپاره، که~$n \neq m$
وي، نو $Q \Proves \eq/[\num n][\num m]$۔

په بنسټيزه قضيه کښې~$n = 0$ دے۔ که~$m$ له~$0$ سره مساوي نۀ وي، نو
د کوم طبيعي عدد~$k$ دپاره~$m = k + 1$۔ يو داسې بديهي اصل لرو چې
$\lforall[x][\eq/[0][x']]$ وايي۔ د کمیت ايښوونکي د يوه بديهي اصل له
مخې، د~$x$ پر ځاے د~$\num k$ په ايښودلو، دا پايله اخلو:
$\eq/[0][\num k']$۔ خو~$\num k'$ هماغه~$\num m$ دے۔

د استقرا په پړاو کښې فرضولے شو چې ادعا د~$n$ دپاره سمه ده، او
$n+1$ ګورو۔ $m$ هر طبيعي عدد ونيسئ۔ دوه امکانونه شته: يا~$m = 0$
دے يا د کوم~$k$ دپاره~$m = k+1$۔ لومړۍ قضيه لکه پورته حل کېږي۔ په
دويمه قضيه کښې فرض کړئ~$n+1 \neq k+1$۔ بيا~$n \neq k$۔ د~$n$
دپاره د استقرايي فرضيې له مخې
$\Th{Q} \Proves \eq/[\num n][\num k]$۔ داسې بديهي اصل لرو چې
$\lforall[x][\lforall[y][\eq[x'][y'] \lif \eq[x][y]]]$ وايي۔ د
کمیت ايښوونکي د يوه بديهي اصل په کارولو، لرو:
$\eq[\num n'][\num k'] \lif \eq[\num n][\num k]$۔ د قضيو د منطق په
کارولو، په~$\Th{Q}$ کښې دا پايله اخلو:
$\eq/[\num n][\num k] \lif \eq/[\num n'][\num k']$۔ د وضع مقدم په
کارولو، $\eq/[\num n'][\num k']$ اخلو، چې هماغه غوښتنه ده، ځکه
$\num k'$، $\num m$ دے۔
\end{proof}

\begin{explain}
پام وکړئ چې لمه ډېره خبره نۀ کوي: په اصل کښې وايي چې~$\Th{Q}$
ثابتولے شي چې بېلابېلې عددنښې بېلابېل څيزونه ښيي۔ د بېلګې په
توګه،~$\Th{Q}$، $0'' \neq 0'''$ ثابتوي۔ خو دا ښودل چې خبره په عام
ډول سمه ده څه پاملرنه غواړي۔ دا هم په پام کښې ونيسئ چې که څه هم
استقرا کاروو، خو دا د~$\Th{Q}$ \emph{بهر} استقرا ده۔
\end{explain}

\begin{proof}[د \olref{prop:rep-id} ثبوت]
که~$n = m$ وي، نو~$\num{n}$ او~$\num{m}$ يو شان ترمونه دي، او
$\Char{=}(n, m) = 1$۔ خو $\Th{Q} \Proves (\eq[\num{n}][\num{m}] \land
\eq[\num{1}][\num{1}])$، نو~$!A_=(\num{n}, \num{m},
\num{1})$ ثابتوي۔ که~$n \neq m$ وي، نو~$\Char=(n, m) = 0$۔ د
\olref{lem:q-proves-neq} له مخې،
$\Th{Q} \Proves \eq/[\num{n}][\num{m}]$، او نو دا هم ثابتوي:
$(\eq/[\num{n}][\num{m}] \land \Obj 0 = \Obj 0)$۔ له دې امله
$\Th{Q} \Proves !A_=(\num{n}, \num{m}, \num{0})$۔

د دويمې برخې دپاره هم دوه قضيې لرو۔ که~$n = m$ وي، بايد وښيو:
$\Th{Q} \Proves \lforall[y][(!A_=(\num{n}, \num{m}, y)
  \lif \eq[y][\num{1}])]$۔ په غيرصوري ډول استدلال کوو، فرض کړئ
$!A_=(\num{n}, \num{m}, y)$، يعنې
\[
(\eq[\num{n}][\num{n}] \land \eq[y][\num{1}]) \lor
(\eq/[\num{n}][\num{n}] \land \eq[y][\num{0}])
\]
کيڼ منفصل د منطق له مخې~$\eq[y][\num{1}]$ لازموي؛ ښے منفصل له
$\eq[\num{n}][\num{n}]$ سره تناقض لري، چې په منطق ثابتېدونکے دے۔

بل پلو، فرض کړئ~$n \neq m$۔ بيا~$!A_=(\num{n},
\num{m}, y)$ دا دے:
\[
(\eq[\num{n}][\num{m}] \land \eq[y][\num{1}]) \lor
(\eq/[\num{n}][\num{m}] \land \eq[y][\num{0}])
\]
دلته کيڼ منفصل له~$\eq/[\num{n}][\num{m}]$ سره تناقض لري، چې د
\olref{lem:q-proves-neq} له مخې په~$\Th{Q}$ کښې ثابتېدونکے دے؛ ښے
منفصل~$\eq[y][\num{0}]$ لازموي۔
\end{proof}

\begin{prop}
\ollabel{prop:rep-add}
د جمع تابعه~$\Add(x_0, x_1) = x_0+x_1$ په~$\Th{Q}$ کښې د دې په
وسيله تمثيلېږي:
\[
  !A_{\Add}(x_0, x_1, y) \ident \eq[y][(x_0 + x_1)].
\]
\end{prop}

\begin{lem}
\ollabel{lem:q-proves-add}
$\Th{Q} \Proves \eq[(\num{n} + \num{m})][\num{n+m}]$
\end{lem}

\begin{proof}
دا پر~$m$ په استقرا ثابتوو۔ که~$m = 0$ وي، ادعا دا ده:
$\Th{Q} \Proves \eq[(\num{n} + \Obj 0)][\num{n}]$۔ دا له بديهي
اصل~$!Q_4$ څخه راځي۔ اوس د~$m$ دپاره ادعا فرض کړئ؛ راځئ د~$m+1$
دپاره ادعا ثابته کړو، يعنې وښيو:
$\Th{Q} \Proves \eq[(\num{n} +
  \num{m+1})][\num{n+m+1}]$۔ پام وکړئ چې~$\num{m+1}$ هماغه
$\num{m}'$ دے، او~$\num{n+m+1}$ هماغه~$\num{n+m}'$ دے۔ د بديهي
اصل~$!Q_5$ له مخې، $\Th{Q} \Proves
\eq[(\num{n} + \num{m}')][(\num{n}+\num{m})']$۔ د استقرايي فرضيې
له مخې، $\Th{Q} \Proves
\eq[(\num{n} + \num{m})][\num{n+m}]$۔ نو
$\Th{Q} \Proves \eq[(\num{n} + \num{m}')][\num{n+m}']$۔
\end{proof}

\begin{proof}[د \olref{prop:rep-add} ثبوت]
هغه !!{formula}~$!A_\Add(x_0, x_1, y)$ چې~$\Add$ تمثيلوي،
$\eq[y][(x_0 + x_1)]$ ده۔ لومړے ښيو چې که~$\Add(n, m) = k$ وي،
نو $\Th{Q} \Proves !A_\Add(\num{n}, \num{m}, \num{k})$؛ يعنې
$\Th{Q} \Proves \eq[\num{k}][(\num{n} + \num{m})]$۔ خو ځکه چې
$k = n + m$، نو~$\num{k}$ هماغه~$\num{n+m}$ دے، او په
\olref{lem:q-proves-add} کښې مو ښودلې ده چې
$\Th{Q} \Proves \eq[(\num{n} + \num{m})][\num{n+m}]$۔

دا هم بايد وښيو چې که~$\Add(n, m) = k$ وي، نو
\[
\Th{Q} \Proves \lforall[y][(!A_\Add(\num{n}, \num{m}, y) \lif
  \eq[y][\num{k}])].
\]
فرض کړئ~$\eq[(\num{n} + \num{m})][y]$ لرو۔ ځکه چې
\[
\Th{Q} \Proves \eq[(\num{n}+\num{m})][\num{n+m}],
\]
کيڼ اړخ په~$\num{n+m}$ بدلولے شو او د خپل سري~$y$ دپاره
$\eq[\num{n+m}][y]$ تر لاسه کوو۔
\end{proof}

\begin{prop}
\ollabel{prop:rep-mult}
د ضرب تابعه~$\Mult(x_0, x_1) = x_0 \cdot x_1$ په~$\Th{Q}$ کښې د
دې په وسيله تمثيلېږي:
\[
  !A_{\Mult}(x_0, x_1, y) \ident y = (x_0 \times x_1).
\]
\end{prop}

\begin{proof} تمرین۔ \end{proof}

\begin{lem}
\ollabel{lem:q-proves-mult}
$\Th{Q} \Proves \eq[(\num{n} \times \num{m})][\num{n \cdot m}]$
\end{lem}

\begin{proof} تمرین۔ \end{proof}

\begin{prob}
\olref[inc][req][bre]{lem:q-proves-mult} ثابته کړئ۔
\end{prob}

\begin{prob}
له \olref[inc][req][bre]{lem:q-proves-mult} څخه د
\olref[inc][req][bre]{prop:rep-mult} په ثبوت کښې کار واخلئ۔
\end{prob}

\begin{explain}
  په ياد ولرئ چې د حساب د ژبې د تابعې نښې دپاره~$\times$، او پر
  عددونو د عادي ضرب د عمل دپاره~$\cdot$ کاروو۔ نو~$\cdot$ د عددونو
  د تعبيرونو ترمنځ راتلے شي (لکه په~$m \cdot n$ کښې)، خو~$\times$
  يوازې د حساب د ژبې د ترمونو ترمنځ راځي (لکه په
  $(\num{m} \times \num{n})$ کښې)۔ تر دې هم مغشوشوونکې دا ده چې
  $+$ هم د !!{function} او هم د جمع د عمل دپاره کارېږي۔ کله چې د
  ترمونو ترمنځ راځي---لکه په~$(\num{n} + \num{m})$ کښې---دا د حساب
  د ژبې~$2$-ځايه !!{function} ده، او کله چې د عددونو ترمنځ راځي---
  لکه په~$n+m$ کښې---دا د جمع عمل دے۔ په دې کښې د~$\num{n+m}$
  قضيه هم شامله ده: دا له عدد~$n+m$ سره سمون لرونکې معياري عددنښه
  ده۔
\end{explain}

\end{document}
